

%%Changes by Pierre begun on
%%February 10, 2003 
%% AMS-TeX file
%% original begun in July, 2000
%% Final corrections: February 20, 2003

\documentstyle{amsppt}
\magnification=1200
\loadeufm
%\input epsf.tex

\font\dotless=cmr10 %for the roman i or j to be
                    %used with accents on top.
                    %(\dotless\char'020=i)
                    %(\dotless\char'021=j)
\font\itdotless=cmti10
\font\mathdotless=cmmi10
\def\itumi{{\"{\itdotless\char'020}}}
\def\itumj{{\"{\itdotless\char'021}}}
\def\umi{{\"{\dotless\char'020}}}
\def\umj{{\"{\dotless\char'021}}}
\def\hati{{\^{\dotless\char'020}}}
\def\Hati{\hbox{$\hat{\imath}$}}
\def\Hatj{\hbox{$\hat{\jmath}$}}
%\font\smaller=cmr5
\font\boldtitlefont=cmssdc10 scaled\magstep2
\font\authorfont=cmss10 scaled\magstep1
\font\sectionfont=cmss10 scaled\magstep1
\font\sans=cmss10


\NoRunningHeads
\pagewidth{6.5 true in}
\pageheight{8.9 true in}
\loadeusm

\catcode`\@=11
\def\logo@{}
\catcode`\@=13


\newif\iffirstul

\def\ullspecial#1{\ifx\empty#1 \let\next=\relax \else
    \iffirstul\firstulfalse\else%
    \discretionary{}{}{\hbox{\ }}\fi%
$\underline{\vphantom{y}%
  \hbox{#1}}$\let\next=\ullspecial\fi 
 \next}

\def\ul#1{\firstultrue\ullspecial#1 \empty }
%
\def\undertext#1{$\underline{\vphantom{y}\hbox{#1}}$}
\def\nspace{\lineskip=1pt\baselineskip=12pt%
     \lineskiplimit=0pt}
\def\dspace{\lineskip=2pt\baselineskip=18pt%
     \lineskiplimit=0pt}

%\def\wedgeop{\operatornamewithlimits{\wedge}\limits}
\def\smbul{{\sssize\bullet}}
\def\oplusop{\operatornamewithlimits{\oplus}\limits}
\def\otimesop{\operatornamewithlimits{\otimes}
     \limits}
%\def\Piop{\operatornamewithlimits{\Pi}\limits}
\def\colim{\operatornamewithlimits{\text{\rm colim}}}
\def\w{{\mathchoice{\,{\scriptstyle\wedge}\,}
{{\scriptstyle\wedge}}{{\scriptscriptstyle\wedge}}%
     {{\scriptscriptstyle\wedge}}}}
\def\Le{{\mathchoice{\,{\scriptstyle\le}\,}
{\,{\scriptstyle\le}\,}{\,{\scriptscriptstyle\le}\,}%
{\,{\scriptscriptstyle\le}\,}}}
\def\Ge{{\mathchoice{\,{\scriptstyle\ge}\,}
{\,{\scriptstyle\ge}\,}{\,{\scriptscriptstyle\ge}\,}%
{\,{\scriptscriptstyle\ge}\,}}}
%\def\doubleheaddownarrow#1{\hbox{$\Big\downarrow%
%     \kern-5.96pt\lower2pt\hbox{$\downarrow$}$}%
%     \rlap{$\vcenter{$\kern-18pt\scriptstyle#1$}$}}
%\def\rmapdown#1{\Big\downarrow\kern-1.0pt\vcenter{
%     \hbox{$\scriptstyle#1$}}}
\def\tworightarrows#1#2{\lower2.0pt\hbox{\,\,\vbox{\baselineskip=0pt
    \hbox{\lower6.0pt\hbox to #1pt{\rightarrowfill}}\kern-6pt
    \hbox to #2pt{\rightarrowfill}}\,\,}}
\def\mapright#1{\smash{\mathop{\,\longrightarrow\,}%
  \limits^{#1}}}
\def\longmapright#1#2{\,\smash{\mathop
     {\hbox to #1pt{\rightarrowfill}}\limits^{#2}}\,}
\def\longmaprightsub#1#2{\,\smash{\mathop
     {\hbox to #1pt{\rightarrowfill}}\limits_{#2}}\,}
\def\simover{\smash{\mathop{\,\longrightarrow\,}
  \limits^{\lower1.5pt\hbox{$\scriptstyle\sim$}}}}
\def\longhook{\lhook\!\longrightarrow}
\def\Righthook#1{\scriptscriptstyle%
  \subset\kern-4.00pt\lower2.85pt
     \hbox to #1pt{\rightarrowfill}}
%\Righthook{30}, for instance
\def\Heading#1#2{\bigskip\par\noindent{\bf#1}
     \medskip\parindent=#2pt\nspace}
\def\Ref[#1]{\par\smallskip\hang\indent\llap{\hbox to\parindent
     {[#1]\hfil\enspace}}\ignorespaces}
\def\ref#1{\par\smallskip\hang\indent\llap{\hbox to\parindent
     {#1\hfil\enspace}}\ignorespaces}

%\def\Cech{\hbox{\v{C}ech}}
%\def\Ccheck{\hbox{\v{C}}}
%\def\(({(\!(}   \def\[[{[\![}
%\def\)){)\!)}   \def\]]{]\!]}
%\def\wdot{\Dot{w}}
%\def\sdot{\Dot{s}}
%\def\xdot{\Dot{x}}
%\def\emptydot{\Dot{\phantom{L}}}
%\def\righthookdown{\raise8.5pt\hbox{$\ssize\cap$}
%     \kern-7.80pt \big\downarrow}
%\def\rhookmapdown#1{\rightcapdownarrow\kern-1.0pt
%   \vcenter{\hbox{$\scriptstyle#1$}}}
\def\Item#1{\par\smallskip\hang\indent
  \llap{\hbox to\parindent{#1\hfill\enspace}}
     \ignorespaces}
\def\Amalg{\,\shortparallel\kern-10.15 true pt
\hbox{\lower.5 true pt\hbox to 4.0 true pt{\hrulefill}}\,\,}
%\def\circled-one{\text{O}\kern-5.5pt\raise2pt%
%     \hbox{$\sssize 1$}}
%\font\smll=cmr5
%\def\Small-circled-one{\hbox{$\tsize\text{O}$}%
%     \kern-5.25 pt\raise1.25pt\hbox{\smll 1}}
%\def\rightcapdownarrow{\raise9pt\hbox{$\ssize\cap$}
%\kern-7.75pt\Big\downarrow}
%\def\vrulesub#1{\hbox{\,\vrule height7pt depth5pt\,}_{#1}}
\def\upvee{^{\sssize\vee}}
\def\uphat{^{\sssize\wedge}}
\def\CUP{{\,\ssize \cup\,}}
\def\CAP{{\,\ssize \cap\,}}
\def\Coprod{{\ssize\coprod}}
%\def\restprod{\raise1.5pt\hbox{${\ssize\prod}
%  \kern-8pt{\ssize\coprod}$}}
%\def\ThreeArrows{\mathop{\vtop{\ialign{##\crcr
%\hfil$\sssize\longrightarrow$\hfil\crcr
%\noalign{\nointerlineskip\vskip.95pt}%
%  $\sssize\longrightarrow$\crcr
%\noalign{\nointerlineskip\vskip.95pt}%
%  $\sssize\longrightarrow$\crcr
%\noalign{\nointerlineskip}\crcr}}}}

%\def\FourArrows{\mathop{\vtop{\ialign{##\crcr
%\hfil$\sssize\longrightarrow$\hfil\crcr
%\noalign{\nointerlineskip\vskip.95pt}
%$\sssize\longrightarrow$\crcr
%\noalign{\nointerlineskip\vskip.95pt}
%$\sssize\longrightarrow$\crcr
%\noalign{\nointerlineskip\vskip.95pt}
%$\sssize\longrightarrow$\crcr
%\noalign{\nointerlineskip}\crcr}}}}

\redefine\Supset{\raise1.75pt
     \hbox{$\,\,\ssize\supset\,\,$}}
\redefine\Subset{\raise1.75pt
     \hbox{$\,\,\ssize\subset\,\,$}}
\redefine\mod{\,\text{\rm mod\,}}
\redefine\Cup{\raise.80pt
     \hbox{$\,\ssize\cup\,$}}



\def\DR{\text{\rm DR}}
\def\Gal{\text{\rm Gal}}
\def\Ext{\text{\rm Ext}}
\def\End{\text{\rm End}}
\def\ab{\text{\rm ab}}
\def\ad{\text{\rm ad}}
\def\comp{\text{\rm comp}}
\def\Dec{\text{\rm Dec}}
\def\distr{\text{\rm distr}}
\def\est{\text{\rm est}}
\def\Spec{\text{\rm Spec}}
\def\Gr{\text{\rm Gr}} 
\def\GR{\text{\rm gr}}
\def\Sym{\text{\rm Sym}} 
\def\GL{\text{\rm GL}} 
\def\Hom{\text{\rm Hom}} 
\def\HOM{{\scrH}\text{\rm om}}
\def\image{\text{\rm image}}
\def\Ind{\text{\rm Ind}}
\def\Id{\text{\rm Id}}
\def\Lie{\text{\rm Lie}}
\def\Libre{\text{\rm Libre}}
\def\Lib{\text{\rm Lib}}
\def\Libhat{\Lib^{\uphat}}
\def\rang{\text{\rm rang}}
\def\ieme{\text{\rm i\`eme}}
\def\iemes{\text{\rm i\`emes}}
\def\im{\text{\rm Im}}
\def\ev{\text{\rm ev}}
\def\et{\text{\rm et}}
\def\Fil{\text{\rm Fil}}
\def\real{\text{\rm real}}
\def\mot{\text{\rm mot}}
\def\Mot{\text{\rm Mot}}
\def\MAT{\text{\rm MAT}}
\def\MTM{\text{\rm MTM}}
\def\MPTM{\text{\rm MPTM}}
\def\MT{\text{\rm MT}}
\def\DM{\text{\rm DM}}
\def\DMT{\text{\rm DMT}}
\def\SmCor{\text{\rm SmCor}}
%\def\Point{\text{\rm Point}}
\def\Sing{\text{\rm Sing}}
\def\torsion{\text{\rm torsion}}
\def\pr{\text{\rm pr}}
\def\Ker{\text{\rm Ker}}
\def\kar{\text{\rm kar}}
\def\Aut{\text{\rm Aut}} 
\def\Rep{\text{\rm Rep}}
\def\Pet{\text{\rm Pet}}
\def\dch{\text{\rm dch}}
\def\ch{\text{\rm ch}}
\def\coeff{\text{\rm coeff}}

\def\eps{{\varepsilon}}

\def\abar{\bar{a}} 
\def\dbar{\bar{d}} 
\def\gbar{\bar{g}}
\def\kbar{\bar{k}} 
\def\Kbar{\bar{K}}
\def\Xbar{\bar{X}}
\def\xbar{\bar{x}}
\def\Ybar{\bar{Y}}
\def\ybar{\bar{y}} 
\def\zbar{\bar{z}}
\def\Zbar{\bar{Z}}
\def\Gammabar{\bar{\Gamma}}
\def\sigmabar{\bar{\sigma}}
\def\zetabar{{\bar{\zeta}}}
\def\dbQbar{\bar{\dbQ}}
\def\underbark{\hbox{\b{$k$}}}
\def\underbart{\hbox{\b{$t$}}}

\def\khat{\Hat{k}}
\def\xhat{\Hat{x}}

\def\bfs{\bold{s}}
\def\bfv{\bold{v}}
\def\bfQ{\bold{Q}}

\def\Etil{\widetilde{E}}
\def\Mtil{\widetilde{M}}
%\def\uptilde{^{\raise3pt\hbox{$\sssize\sim$}}}
\def\uptil{\kern-1pt\tilde{\phantom{a}}}
\def\upcheck{^{\sssize\vee}}
\def\Lamuptil{\Lambda\raise1pt\hbox{$\!\uptil$}}

\def\dbA{{\Bbb A}} 
\def\dbC{{\Bbb C}} 
\def\dbG{{\Bbb G}} 
\def\dbH{{\Bbb H}} 
\def\dbR{{\Bbb R}} 
\def\dbP{{\Bbb P}}
\def\dbQ{{\Bbb Q}}
\def\dbZ{{\Bbb Z}}

\def\scr#1{{\fam\eusmfam\relax#1}}

\def\scrA{{\scr A}} 
\def\scrC{{\scr C}} 
\def\scrD{{\scr D}}
\def\scrE{{\scr E}}
\def\scrF{{\scr F}}
\def\scrH{{\scr H}}
\def\scrK{{\scr K}}
\def\scrI{{\scr I}}
\def\scrL{{\scr L}}
\def\scrM{{\scr M}}
\def\scrO{{\scr O}}
\def\scrR{{\scr R}}
\def\scrS{{\scr S}}
\def\scrT{{\scr T}}
\def\scrU{{\scr U}}

\def\gr#1{{\fam\eufmfam\relax#1}}

%Euler Fraktur letters (German)
\def\gra{{\gr a}}
\def\gru{{\gr u}} 
\def\grg{{\gr g}} 
\def\grl{{\gr l}}
\def\grs{{\gr s}}
\def\grQ{{\gr Q}}
\def\grZ{{\gr Z}}
\def\grgl{{\grg\grl}}

%\NoBlackBoxes
\parindent=20pt
\frenchspacing
\footline={\hfill}

\topmatter
\title\nofrills
{\boldtitlefont Groupes fondamentaux motiviques
de Tate mixte}
\endtitle
\bigskip
\author
{\authorfont P. Deligne et A.~B.~Goncharov$^*$}
\endauthor
\endtopmatter
\footnote""{$^*$Research partially supported by NSF
Grant DMS-0099390}

\NoBlackBoxes
\document

{\narrower{\narrower{\narrower{
\bigskip\bigskip\bigskip\noindent
{\sans R\'esum\'e.}\enspace
Nous d\'efinissons la cat\'egorie des motifs de
Tate mixte sur l'anneau des $S$-entiers d'un 
corps de nombres, et le groupe fondamental
motivique (rendu unipotent) d'une vari\'et\'e
unirationnelle sur un corps de nombres.
Nous consid\'erons plus en d\'etail le groupe
fondamental motivique de la droite projective
moins $0$, $\infty$ et les racines $N$-i\`emes
de l'unit\'e.\bigskip}}}}

\bigskip\bigskip
{\narrower{\narrower{\narrower{\bigskip\noindent
{\sans Abstract.}\enspace
We define the category of mixed Tate motives
over the ring of $S$-integers of a number
field.
We define the motivic fundamental group (made
unipotent) of a unirational variety over a
number field.
We apply this to the study of the motivic
fundamental group of the projective line
punctured at zero, infinity and all $N$-th
roots of unity.\bigskip}}}}


\newpage
\footline={\hss\folio\hss}
\pageno=-1

\def\leaderfill{\leaders\hbox to 1em{\hss.\hss}\hfill}

\topmatter
\toc\nofrills{\sectionfont Table of Contents}
\bigskip\bigskip
%\openup.5\baselineskip
\widestnumber\head{2a.}
\head 0. Introduction\page{1}\endhead
\bigskip
\head 1. Motifs de Tate mixte sur l'anneau
     des $S$-entiers de $k$ \page{3}\endhead
\bigskip
\head 2. Le point de vue tannakien \page{13}\endhead
\bigskip
\head 3.  Cohomologie et groupe fondamental
  \page{25}\endhead
\bigskip
\head 4. Le groupe fondamental unipotent
motivique d'une\hfill\break
vari\'et\'e rationnelle\page{40}\endhead
\bigskip
\head 5. Le groupe fondamental motivique du
compl\'ement,\hfill\break
      dans $\dbP^1$, de $0$, $\infty$ et
     $\mu_N$\page{50}\endhead
\bigskip
\head 6. Profondeur\page{67}\endhead
\bigskip
\specialhead{} Appendice. Rappels sur les
groupes alg\'ebriques 
   unipotents \page{76}\endspecialhead
\bigskip
\specialhead{} Bibliographie\page{82}\endspecialhead
\endtoc

\endtopmatter

\input introccc.tex 
\input Sec1ccc.tex
\input Sec2ccc.tex
\input Sec3ccc.tex
\input Sec4ccc.tex
\input Sec5ccc.tex
\input Sec6ccc.tex
\input Appenccc.tex
\input Refs.tex

\enddocument




